\subsection{里斯自同态的扰动}

引理 1. --- 设 E 是巴拿赫空间。设 p 和 q 是 E 上的连续投影，满足 \(\|q - p\| < 1\) 。则 p 诱导从 \(\operatorname{Im}(q)\) 到 \(\operatorname{Im}(p)\) 上的一个同构，而 \(1_{E} - p\) 诱导从 \(\operatorname{Ker}(q)\) 到 \(\operatorname{Ker}(p)\) 上的一个同构。

考察连续线性映射 \(u\colon x\mapsto p(x)\) （从 \(\operatorname{Im}(q)\) 到 \(\operatorname{Im}(p)\) 中）与 \(v\colon y\mapsto q(y)\) （从 \(\operatorname{Im}(p)\) 到 \(\operatorname{Im}(q)\) 中）。对每个 \(x\in\operatorname{Im}(q)\) ，有 \(x=q(x)\) ，由此

\[
\begin{array}{c} (q - p) ^ {2} (x) = q ^ {2} (x) - q (p (x)) - p (q (x)) + p ^ {2} (x) \\ = x - q (p (x)) = x - v (u (x)). \end{array}
\]

由此可得 \(\|1_{E}-v\circ u\|\leqslant\|q-p\|^{2}<1\) 。由 I, p. 22 的推论 1，\(v\circ u\) 是 \(\operatorname{Im}(q)\) 的一个自同构。类似地可证 \(u\circ v\) 是 \(\operatorname{Im}(p)\) 的一个自同构。这蕴含 u 是从 \(\operatorname{Im}(q)\) 到 \(\operatorname{Im}(p)\) 上的一个同构，从而得出引理的第一个断言。第二个断言由此通过把 p 和 q 分别换成 \(1_{E}-p\) 和 \(1_{E}-q\) 而得到。

定理 2. --- 设 E 是巴拿赫空间。E 的里斯自同态所成的集 \(\mathcal{R}(\mathrm{E})\) 在 \(\mathcal{L}(\mathrm{E})\) 中是开的。把 \(\mathcal{R}(\mathrm{E})\) 的元素映为其幂零空间的维数的映射在 \(\mathcal{L}(\mathrm{E})\) 上是上半连续的。

定理由下面更精确的命题得出：

命题 5. --- 设 E 是巴拿赫空间，\(u_{0}\) 是 E 的一个里斯自同态。设 N（相应地，I）是 \(u_{0}\) 的幂零空间（相应地，余幂零空间）。用 d 表示 N 的维数。存在 \(u_{0}\) 在 \(\mathcal{L}(\mathrm{E})\) 中的一个开邻域 U 和一个解析映射 \(\pi\colon\mathrm{U}\to\mathcal{L}(\mathrm{E})\) ，使得

\begin{enumerate}
\def\labelenumi{(\roman{enumi})}
\item
  自同态 \(\pi(u_0)\) 是以 N 为像、以 I 为核的投影；
\item
  对每个 \(u \in U\) ，线性映射 \(\pi(u)\) 是秩为 d 的投影，它属于 u 的双交换子，特别地与 u 交换。此外，u 诱导 \(\operatorname{Ker}(\pi(u))\) 的一个自同构；
\item
  U 的每个元素都是里斯自同态，其幂零空间的维数 \(\leqslant d\) 。
\end{enumerate}

若 K = C，用 \(\mathrm{Sp}(u_{0})\) 表示 \(u_{0}\) 关于代数 \(\mathcal{L}(\mathrm{E})\) 的谱。当 \(K = R\) 时，用 \(\mathrm{Sp}(u_{0})\) 表示复谱 \(\mathrm{Sp}_{\mathcal{L}(\mathrm{E}_{(\mathbf{C})})}((u_{0})_{(\mathbf{C})})\) ，即 \(u_{0}\) 的复谱（I, p. 85, 第 13 目）。由 III, p. 54 的命题 14，点 0 在 \(\{0\} \cup \mathrm{Sp}(u_{0})\) 中是孤立的。设 r \textgreater{} 0 是实数，使得 \(\mathrm{Sp}(u_{0}) - \{0\}\) 的每个元素的模都 \textgreater{} r。用 V 表示 C 中由绝对值 ≠ r 的复数所组成的开集；用 f 表示 V 上如下定义的全纯函数：\(f(z) = 0\) （若 \(|z| > r\) ），\(f(z) = 1\) （若 \(|z| < r\) ）。若 K = C（相应地，K = R），用 U' 表示 \(\mathcal{L}(\mathrm{E})\) 中谱（相应地，复谱）包含于 V 的元素 u 所成的集。集 U' 是 \(u_{0}\) 在 \(\mathcal{L}(\mathrm{E})\) 中的一个开邻域，且映射 \(u \mapsto f(u)\) （从 U' 到 \(\mathcal{L}(\mathrm{E})\) 中）是全纯的（I, p. 76, 命题 10 及 I, p. 85, 第 13 目）。

设 \(u \in \mathrm{U}'\) 。自同态 \(f(u)\) 是谱投影 \(e_0(u)\) ；它与 \(u\) 交换，且 \(u\) 通过过渡到子空间诱导 \(\operatorname{Ker}(e_0(u))\) 的一个自同构（参见 III, p. 53, 第 6 目）。

投影 \(e_{0}(u_{0})\) 以 N 为像、以 I 为核（III, p. 54, 命题 14）；它的秩是 d。由引理 1，由 \(u \in U'\) 中使 \(e_{0}(u)\) 的秩为 d 的元素所成的集 U 是 \(u_{0}\) 的一个开邻域。集 U 与映射 \(\pi: u \mapsto e_{0}(u)\) （从 U 到 \(\mathcal{L}(\mathrm{E})\) 中）满足命题的条件 (i) 与 (ii)。

设 \(u \in U\) 。由于 u 诱导 \(\operatorname{Ker}(e_{0}(u))\) （它是 E 的一个闭的有限余维数向量子空间）的一个自同构，自同态 u 是 E 的里斯自同态（III, p. 48, 命题 8）；u 的幂零空间包含于 \(\pi(u) = e_{0}(u)\) 的像中，且其维数 \(\leqslant d\) 。

